\documentclass[a4paper]{article}

\usepackage{graphicx}
\usepackage{amsmath,amssymb}
\usepackage{minted}
\usepackage{multirow}
\usepackage{float}
\usepackage{algorithm}
\usepackage{algpseudocode}
\usepackage{todonotes}
\usepackage{forest}
\usepackage{xstring}
\usepackage{xcolor}
\usepackage[most]{tcolorbox}
\usepackage{xparse}
\usepackage{natbib}
\usepackage{adjustbox}
\usepackage[active,tightpage]{preview}

\PreviewBorder=4pt

\usetikzlibrary{graphs,graphdrawing,arrows.meta,positioning,calc,fit,shapes.geometric}
\usegdlibrary{layered}

\input{/nd_language_ai_project/src/nd_language_ai/formal/computer/programming/tex/latex/code_clips/external_graphviz.tex}
\input{/nd_language_ai_project/src/nd_language_ai/formal/computer/programming/tex/latex/part/control_sequence/kind/semtag/semtag.tex}
\input{/nd_language_ai_project/src/nd_language_ai/formal/computer/programming/tex/latex/code_clips/seehere.tex}
\input{/nd_language_ai_project/src/nd_language_ai/formal/computer/programming/tex/latex/code_clips/deep_structure.tex}

\title{Large Visual Structure Initial-load Fix}
\date{2026-09-28}
\author{Mohammad Rahmani}

\begin{document}
   \maketitle

   \begin{preview}
      \section{Cause}
      Restoring a virtualized large LaTeX document directly into Visual mode
      initialized the mode-isolation cache from the materialized Source
      \texttt{QTextDocument}. That document contains only the roughly
      40,000-character Source window, not the complete canonical buffer.
      A deferred content emission could therefore send only that window to
      Structure. For \texttt{evil.tex}, the first Source window ended near
      line 154 and produced exactly the first seven Structure nodes through
      the subsection at line 152.

      \section{Canonical-source ownership}
      Visual mode now obtains and emits content through the editor's canonical
      source provider when one exists. Large editors therefore keep the full
      canonical buffer as the sole source for Structure, saving, Preview, and
      delayed content notifications while Source remains virtualized.

      \section{Presentation-only document changes}
      Detaching Source syntax highlighting during the Visual transition can
      produce a same-length \texttt{QTextDocument.contentsChange} event even
      though no text changed. Large Source synchronization now identifies that
      presentation-only event and does not apply it to the canonical buffer.
      This prevents a mode switch from changing source text or marking the
      document as modified.

      \section{Verification}
      A regression test restores a source larger than the virtualization
      threshold directly into Visual mode, verifies that the materialized
      Source remains smaller than the canonical document, forces the deferred
      content path, and confirms that Structure receives every section.
      The real \texttt{evil.tex} startup path was also exercised: the
      canonical source remained unchanged and unmodified, and the first
      Structure parse produced 102 nodes through the final section.
   \end{preview}

   \bibliographystyle{plainnat}
   \bibliography{/home/donkarlo/Dropbox/repo/nd_shared_research/bibliography/bibliography.bib}
\end{document}
